% chapters_parte2/08_multicast.tex

\chapter{\eh{mega} Multicast}

Imagina transmitir un partido de fútbol a 1 millón de espectadores. Con \textbf{unicast},
el servidor enviaría 1 millón de copias del mismo stream: imposible. Con
\textbf{broadcast}, llegaría a todos los dispositivos de la red, incluso a quienes no
quieren verlo. \textbf{Multicast} resuelve esto elegantemente: el servidor envía
\emph{una sola copia} que la red replica solo hacia quienes la solicitaron.

% ─────────────────────────────────────────────
\section{\eh{bar-chart} Unicast vs Broadcast vs Multicast}

\begin{center}
\begin{tabular}{llll}
\toprule
\textbf{Modelo} & \textbf{Destinatario} & \textbf{Copias en red} & \textbf{Uso} \\
\midrule
Unicast   & Un host específico       & Una por destino         & HTTP, SSH, DNS \\
Broadcast & Todos en el segmento L2  & Una (pero llega a todos) & ARP, DHCP \\
Multicast & Grupo de interesados     & Una (la red la replica) & IPTV, routing protocols \\
Anycast   & El más cercano del grupo & Una                     & DNS raíz, CDN \\
\bottomrule
\end{tabular}
\end{center}

\begin{cuidado}
Multicast tiene reputación de ser complejo de implementar y operar. No es mito:
requiere soporte en cada router del camino, gestión de grupos, y protocolos adicionales
(PIM, IGMP, MSDP). Por eso en la nube pública casi no existe multicast nativo: AWS
y GCP no lo soportan en VPCs. Se reemplaza con overlay multicast (software) o CDN.
\end{cuidado}

% ─────────────────────────────────────────────
\section{\eh{label} Direcciones Multicast}

\begin{concepto}
IPv4 reserva el rango \textbf{224.0.0.0/4} para multicast ($224.x.x.x$ a $239.x.x.x$).
Cada \emph{grupo multicast} tiene una dirección en este rango. Un host se ``une'' al
grupo para recibir tráfico; puede abandonarlo en cualquier momento.
\end{concepto}

\begin{center}
\begin{tabular}{lll}
\toprule
\textbf{Rango} & \textbf{Nombre} & \textbf{Uso} \\
\midrule
224.0.0.0/24   & Link-local       & Protocolos de routing (OSPF: 224.0.0.5/6, HSRP: 224.0.0.2) \\
224.0.1.0/24   & Internetwork     & NTP: 224.0.1.1, uso general \\
232.0.0.0/8    & Source-Specific  & SSM (Source-Specific Multicast) \\
239.0.0.0/8    & Admin-scoped     & Multicast privado dentro de organización \\
\bottomrule
\end{tabular}
\end{center}

\subsection{MAC Multicast}

\begin{concepto}
Cuando un router L2 recibe un paquete multicast IP, lo encapsula en un frame Ethernet
con una MAC de destino especial. El bit I/G del primer byte de la MAC indica multicast.
Para IPs multicast, la MAC se construye así:

\begin{center}
\texttt{01:00:5E:[0][últimos 23 bits de la IP multicast]}
\end{center}

Ej: IP \texttt{224.1.1.1} $\to$ MAC \texttt{01:00:5E:01:01:01}.

Consecuencia: los switches L2 no hacen multicast por defecto --- inundan a todos los
puertos como si fuera broadcast. \textbf{IGMP Snooping} soluciona esto.
\end{concepto}

% ─────────────────────────────────────────────
\section{\eh{ear-of-corn} IGMP --- Gestión de Grupos}

\begin{concepto}
\textbf{IGMP (Internet Group Management Protocol)} es el protocolo entre hosts y el
router de primer salto para gestionar membresía a grupos multicast.
\end{concepto}

\begin{center}
\begin{tabular}{lll}
\toprule
\textbf{Versión} & \textbf{Característica clave} & \textbf{Uso} \\
\midrule
IGMPv1 & Join básico, no hay Leave & Obsoleto \\
IGMPv2 & Leave Group message      & Aún común \\
IGMPv3 & Source filtering (SSM)   & Moderno, requerido para SSM \\
\bottomrule
\end{tabular}
\end{center}

\begin{itemize}
    \item \textbf{IGMP Query:} el router pregunta ``¿alguien quiere grupos multicast?''
          cada 60--125 segundos.
    \item \textbf{IGMP Report:} el host responde con los grupos a los que pertenece.
    \item \textbf{IGMP Leave:} el host notifica que abandona un grupo (IGMPv2+).
    \item \textbf{IGMP Snooping:} los switches L2 ``escuchan'' mensajes IGMP y aprenden
          qué puerto tiene interesados en cada grupo, evitando flooding innecesario.
\end{itemize}

% ─────────────────────────────────────────────
\section{\eh{deciduous-tree} PIM --- Routing Multicast entre Routers}

\begin{concepto}
IGMP maneja el último salto (host $\leftrightarrow$ router). Para que el tráfico
multicast fluya entre routers en la red, se usa \textbf{PIM (Protocol Independent
Multicast)}. ``Protocol Independent'' porque funciona sobre cualquier protocolo de
routing unicast (OSPF, BGP, etc.), usando su tabla de routing para construir árboles
de distribución multicast.
\end{concepto}

\subsection{PIM Dense Mode (PIM-DM)}

\begin{itemize}
    \item Asume que todos los segmentos quieren el tráfico multicast.
    \item \textbf{Flood and prune:} inunda todo, luego los routers sin interesados
          envían mensajes Prune para cortar el árbol.
    \item Simple pero ineficiente. Solo adecuado para redes pequeñas donde la mayoría
          de los segmentos sí tienen receptores.
\end{itemize}

\subsection{PIM Sparse Mode (PIM-SM)}

\begin{concepto}
Modelo opuesto: asume que \emph{pocos} segmentos quieren el tráfico. Los routers
deben suscribirse explícitamente. Usa un \textbf{RP (Rendezvous Point)}: router central
donde fuentes y receptores se ``encuentran''.
\end{concepto}

\begin{center}
\begin{tikzpicture}[
    rtr/.style={draw=cyan!50!white, fill=darkcardgray, text=textlight,
                circle, minimum size=0.9cm, font=\tiny\bfseries},
    rp/.style={draw=yellow!60!white, fill=darkcardviolet, text=textlight,
               circle, minimum size=0.9cm, font=\tiny\bfseries},
    host/.style={draw=gray!50!white, fill=darkcardgray, text=textlight,
                 rectangle, minimum width=1.2cm, minimum height=0.6cm, font=\tiny, align=center},
    arrow/.style={->, thick, draw=gray!60!white},
    joinpath/.style={->, thick, draw=green!60!white, dashed},
    srcpath/.style={->, thick, draw=yellow!60!white},
]
    \node[rp]  (rp) at (3, 4)   {RP};
    \node[rtr] (r1) at (0, 2)   {R1};
    \node[rtr] (r2) at (3, 2)   {R2};
    \node[rtr] (r3) at (6, 2)   {R3};
    \node[host] (src) at (0, 0.5) {Fuente\\224.1.1.1};
    \node[host] (rx1) at (3, 0.5) {Receptor 1};
    \node[host] (rx2) at (6, 0.5) {Receptor 2};
    \draw[arrow] (r1) -- (rp);
    \draw[arrow] (r2) -- (rp);
    \draw[arrow] (r3) -- (rp);
    \draw[srcpath] (src) -- node[font=\tiny, text=yellow!70!white, left] {stream} (r1);
    \draw[joinpath] (rx1) -- node[font=\tiny, text=green!70!white, right] {Join} (r2);
    \draw[joinpath] (rx2) -- node[font=\tiny, text=green!70!white, right] {Join} (r3);
    \draw[srcpath] (rp) -- (r2);
    \draw[srcpath] (rp) -- (r3);
\end{tikzpicture}

{\small Receptores hacen Join hacia el RP. Fuente registra su stream en el RP.
RP distribuye a los routers con receptores. Después puede formarse un Shortest Path Tree
directo fuente $\to$ receptor (SPT switchover).}
\end{center}

\subsection{PIM Source-Specific Multicast (PIM-SSM)}

\begin{concepto}
Variante moderna que elimina el RP. Los receptores se suscriben a un grupo multicast
\emph{y} una fuente específica: \texttt{(S, G)} en lugar de \texttt{(*, G)}. Más simple,
más seguro (no hay fuentes no autorizadas), y usa el rango 232.0.0.0/8.
Requiere IGMPv3 en los hosts.
\end{concepto}

% ─────────────────────────────────────────────
\section{\eh{television} Aplicaciones Reales}

\begin{ejemplo}
\textbf{IPTV (Televisión sobre IP):} operadores como Izzi, Totalplay o Claro TV
distribuyen canales de TV usando multicast. Un canal de HD requiere $\approx$4--8 Mbps.
Con unicast, 10,000 suscriptores viendo el mismo canal = 80 Gbps desde el servidor.
Con multicast = 8 Mbps independientemente del número de suscriptores. El router del
cliente (CPE) hace IGMP Join cuando el usuario cambia de canal.
\end{ejemplo}

\begin{ejemplo}
\textbf{Protocolos de routing:} OSPF usa \texttt{224.0.0.5} (todos los routers OSPF)
y \texttt{224.0.0.6} (DR/BDR). Esto es multicast link-local: el paquete nunca pasa
de ese segmento L2. Sin multicast, OSPF tendría que usar broadcast o unicast punto a punto.
\end{ejemplo}

\begin{moderno}
\emoji{rocket} \textbf{Multicast en datacenters modernos:} VxLAN con BUM traffic
(Broadcast, Unknown-unicast, Multicast) puede usar multicast IP underlay en lugar
de ingress replication. El switch Leaf hace IGMP Join para el grupo multicast del VNI.
Sin embargo, EVPN con ingress replication es hoy más común porque no requiere multicast
en el underlay (muchos clouds no lo soportan). En Kubernetes, \textbf{no hay multicast
nativo en las VPCs de nube pública}: se simula con servicios como AWS SNS (fan-out de
mensajes) o con aplicación-level multicast en frameworks como Hazelcast.
\end{moderno}

\section{\eh{check-mark-button} Resumen del Ecosistema Parte 2}

Al completar estos apuntes has cubierto el stack completo del networking moderno:

\begin{center}
\begin{tabular}{ll}
\toprule
\textbf{Tema} & \textbf{Dónde aplica hoy} \\
\midrule
Wireless empresarial   & Oficinas, campus, IoT industrial \\
Switching avanzado     & Cualquier red LAN empresarial \\
SDN / OpenFlow         & Datacenter, WAN empresarial \\
VxLAN + EVPN           & AWS, Azure, GCP internamente; datacenter privado \\
MPLS + QoS             & ISPs, WAN empresarial, voz sobre IP \\
Kubernetes networking  & Todo cluster en producción \\
Cloud networking       & Toda aplicación en AWS/GCP/Azure \\
Multicast              & IPTV, protocolos de routing, streaming masivo \\
\bottomrule
\end{tabular}
\end{center}
